\chapter{Kimia Supramolekul}\label{ch:b3:supramolecular}

Pada 1962 Charles Pedersen, yang sedang mencari sesuatu yang lain, mengisolasi sebuah
polieter siklik yang melarutkan kalium permanganat dalam benzena dan membuatnya
ungu. Cincin itu membungkus ion kalium dan menyembunyikan muatannya
di balik permukaan hidrokarbon, sambil menyeret permanganat sebagai
ion pengimbang. Kimia di luar molekul, yaitu kimia rakitan-rakitan yang disatukan
oleh gaya yang lebih lemah daripada ikatan kovalen, bermula dari larutan ungu itu.
Bab ini mengukur gaya-gaya itu, menjelaskan mengapa sebagian \omterm{def:b3:supramolecular:supramolecular}{inang} memilih \omterm{def:b3:supramolecular:supramolecular}{tamunya}
dengan begitu baik, memperlihatkan bagaimana tetapan pengikatan diukur, dan ditutup dengan molekul-molekul
yang terulir dan saling mengunci menjadi mesin.

\begin{recall}
Jilid Tahun 1 membahas gaya antarmolekul, ikatan hidrogen, solvasi,
molekul hidrofobik, kompleks, tetapan pembentukan, dan kelat; Jilid Tahun 2
membahas potensial kimia, energi Gibbs, dan pencocokan kuadrat terkecil.
\Cref{ch:b3:advanced-nmr} menjelaskan pertukaran cepat dan lambat serta koalesensi,
\cref{ch:b3:rate-theories} \omterm{thm:b3:rate-theories:eyring}{persamaan Eyring}, dan
\cref{ch:b3:partition-functions} entropi translasi gas.
\end{recall}

\section{Interaksi nonkovalen}

\begin{definition}[Kimia supramolekul]\label{def:b3:supramolecular:supramolecular}
Yang disebut \emph{kimia supramolekul}\index{kimia supramolekul} adalah kimia
rakitan molekul yang disatukan oleh interaksi nonkovalen. Dalam
\emph{kompleks inang--tamu}\index{kompleks inang--tamu}, \emph{inang}\index{inang}
adalah molekul yang lebih besar, dengan rongga atau celah dan situs-situs pengikatan yang mengarah
kepadanya, dan \emph{tamu}\index{tamu} adalah molekul atau ion lebih kecil yang diikatnya.
\end{definition}

Interaksi-interaksinya adalah interaksi pada Jilid Tahun 1, yaitu ion--dipol, dipol--dipol,
ikatan hidrogen, dan gaya London, ditambah beberapa interaksi yang memiliki nama sendiri.

\begin{definition}[Interaksi nonkovalen]\label{def:b3:supramolecular:interactions}
Yang disebut \emph{penumpukan $\pi$}\index{penumpukan $\pi$} adalah tarikan antara cincin aromatik datar
yang saling berhadapan muka, biasanya bergeser atau miring. Yang disebut \emph{interaksi
kation--$\pi$}\index{interaksi kation--$\pi$} adalah tarikan antara kation dan
muka cincin aromatik yang kaya elektron. Yang disebut \emph{ikatan halogen}\index{ikatan halogen}
adalah tarikan antara ujung miskin elektron sebuah atom halogen yang terikat secara kovalen,
di seberang ikatannya, dan sebuah basa Lewis. Yang disebut \emph{efek
hidrofobik}\index{efek hidrofobik} adalah kecenderungan permukaan-permukaan nonpolar untuk
berkumpul di dalam air, yang terutama didorong oleh pelepasan molekul-molekul air yang tertata
di sekelilingnya.
\end{definition}

Interaksi ion--dipol dan kation--$\pi$ adalah yang terkuat, dan sangat bergantung
pada pelarut yang bersaing memperebutkan ion itu; ikatan hidrogen menyusul,
lalu penumpukan $\pi$ dan \omterm{def:b3:supramolecular:interactions}{ikatan halogen}, lalu gaya London, yang lemah satu per satu tetapi
terjumlah atas permukaan kontak yang luas. Di dalam air, \omterm{def:b3:supramolecular:interactions}{efek hidrofobik} sering
mendominasi: \omterm{def:b3:supramolecular:supramolecular}{inang} yang dirancang untuk mengikat \omterm{def:b3:supramolecular:supramolecular}{tamu} dengan ikatan hidrogen dalam kloroform mungkin sama sekali tidak
mengikatnya di dalam air, yang ikatan hidrogennya sendiri ikut bersaing.

\section{Pengenalan molekul}

\begin{definition}[Pengenalan molekul]\label{def:b3:supramolecular:recognition}
Yang disebut \emph{pengenalan molekul}\index{pengenalan molekul} adalah pengikatan
\omterm{def:b3:supramolecular:supramolecular}{tamu} secara selektif oleh \omterm{def:b3:supramolecular:supramolecular}{inang}. Pengenalan itu bergantung pada \emph{komplementaritas}\index{komplementaritas},
yaitu kecocokan ukuran, bentuk, dan situs pengikatan antara \omterm{def:b3:supramolecular:supramolecular}{inang} dan \omterm{def:b3:supramolecular:supramolecular}{tamu}, dan
diperkuat oleh \emph{praorganisasi}\index{praorganisasi}, yaitu \omterm{def:b3:supramolecular:supramolecular}{inang} yang situs-situs
pengikatannya sudah dipertahankan dalam geometri kompleks sebelum pengikatan.
\end{definition}

Tetapan asosiasi $K_a = [\mathrm{HG}]/([\mathrm H][\mathrm G])$ adalah tetapan pembentukan dalam arti
Jilid Tahun 1, dan $\Delta G^\circ = -RT\ln K_a$. \omterm{def:b3:supramolecular:supramolecular}{Inang} yang lentur harus membekukan
konformasinya saat mengikat, dan hal itu mengorbankan entropi; \omterm{def:b3:supramolecular:supramolecular}{inang} yang kaku dan terpraorganisasi telah
membayar harga itu dalam sintesisnya.

\begin{definition}[Efek kelat dan efek makrosiklik]\label{def:b3:supramolecular:macrocyclic}
Yang disebut \emph{efek kelat}\index{efek kelat} adalah kestabilan yang lebih besar sebuah
kompleks dengan ligan polidentat dibandingkan dengan sejumlah ligan
monodentat yang setara dengan atom donor yang sama. Yang disebut \emph{efek
makrosiklik}\index{efek makrosiklik} adalah stabilisasi tambahan ketika
ligan polidentat itu berupa cincin dan bukan rantai terbuka.
\end{definition}

\begin{proposition}[Efek kelat terutama bersifat entropik]\label{prop:b3:supramolecular:chelate-entropy}
Dalam $\mathrm{[M(NH_3)_6]^{2+} + 3\,en \rightleftharpoons [M(en)_3]^{2+} + 6\,NH_3}$, dengan en adalah etana-1,2-diamina, keenam
ikatan logam--nitrogen yang sama terdapat di kedua sisi, dan reaksinya didorong oleh
bertambahnya jumlah molekul bebas.
\end{proposition}

\begin{proof}[Argumen]
Ikatan yang terbentuk dan yang putus sejenis, sehingga $\Delta_rH^\circ$ kecil. Jumlah
molekul yang saling bebas berubah dari empat menjadi tujuh: tiga derajat kebebasan
translasi molekul utuh lebih banyak, yang masing-masing bernilai puluhan joule per kelvin
per mol dalam larutan (\cref{ch:b3:partition-functions}), sehingga $\Delta_rS^\circ > 0$ dan $K > 1$.
Dilihat secara kinetika, begitu satu ujung kelat terikat, ujung yang lain dipertahankan di dekat
logam pada konsentrasi efektif yang tinggi, dan segera terikat kembali jika terlepas.
\end{proof}

\begin{proposition}[Kompensasi entalpi--entropi]\label{prop:b3:supramolecular:compensation}
Dalam serangkaian \omterm{def:b3:supramolecular:supramolecular}{kompleks inang--tamu} yang berkerabat, $\Delta H^\circ$ pengikatan yang lebih negatif
biasanya disertai $\Delta S^\circ$ yang lebih negatif, sehingga $\Delta G^\circ$ berubah lebih sedikit daripada keduanya (sebuah
pengamatan eksperimen).
\end{proposition}

Kompleks yang lebih rapat mengikat lebih kuat tetapi lebih membatasi gerak \omterm{def:b3:supramolecular:supramolecular}{inang} dan
\omterm{def:b3:supramolecular:supramolecular}{tamu}; ikatan hidrogen yang lebih kuat dengan \omterm{def:b3:supramolecular:supramolecular}{tamu} sering berarti ikatan yang lebih lemah dengan
air yang digantikannya. Energi Gibbs, yang menentukan selektivitas, adalah selisih kecil
dua suku yang besar.

\begin{definition}[Inang]\label{def:b3:supramolecular:hosts}
Yang disebut \emph{eter mahkota}\index{eter mahkota} adalah polieter makrosiklik, seperti
18-mahkota-6, $\ce{(CH2CH2O)6}$. Yang disebut \emph{kriptan}\index{kriptan} adalah polieter
bisiklik dengan dua kepala jembatan nitrogen, yang melingkupi kation dalam tiga
dimensi di dalam sebuah \emph{kriptat}\index{kriptat}. Yang disebut
\emph{siklodekstrin}\index{siklodekstrin} adalah oligomer siklik satuan-satuan glukosa,
berbentuk kerucut terpancung dengan rongga hidrofobik dan gugus-gugus hidroksi pada
bibir-bibirnya. Yang disebut \emph{kaliksarena}\index{kaliksarena} adalah makrosiklus berbentuk mangkuk dari
satuan-satuan fenol yang dihubungkan oleh jembatan metilena.
\end{definition}

\begin{omfigure}
\begin{tikzpicture}[font=\small]
  % 18-mahkota-6 dengan K+
  \foreach \k in {0,...,17} {
    \pgfmathsetmacro\ang{90 + 20*\k}
    \pgfmathsetmacro\nxt{110 + 20*\k}
    \draw[black!70] (\ang:1.55) -- (\nxt:1.55);
  }
  \foreach \k in {0,3,...,15} {
    \pgfmathsetmacro\ang{90 + 20*\k}
    \node[fill=white, inner sep=0.5pt, text=omThm] (o\k) at (\ang:1.55) {O};
    \draw[omThm!60, densely dashed] (0,0) -- (\ang:1.32);
  }
  \node[circle, fill=omDef!25, draw=omDef, inner sep=1pt] at (0,0) {$\ce{K+}$};
  \node[font=\scriptsize] at (0,-2.1) {18-mahkota-6 dengan $\ce{K+}$};
  % [2.2.2]kriptan dengan K+
  \begin{scope}[xshift=5.2cm]
    \node (nt) at (0,1.7) {N};
    \node (nb) at (0,-1.7) {N};
    \foreach \x/\n in {-1.25/a, 0/b, 1.25/c} {
      \node[text=omThm, fill=white, inner sep=0.5pt] (oa\n) at (\x,0.55) {O};
      \node[text=omThm, fill=white, inner sep=0.5pt] (ob\n) at (\x,-0.55) {O};
      \draw[black!70] (nt) to[out=-90+40*\x, in=90] (oa\n) -- (ob\n) to[out=-90, in=90-40*\x] (nb);
    }
    \node[circle, fill=omDef!25, draw=omDef, inner sep=1pt] at (0.62,0) {$\ce{K+}$};
    \node[font=\scriptsize] at (0,-2.4) {[2.2.2]kriptan: kriptat $\ce{K+}$};
  \end{scope}
\end{tikzpicture}

{\small Kiri: 18-mahkota-6, cincin dua belas atom karbon (sudut-sudut) dan enam atom
oksigen, dengan ion kalium di pusatnya yang diikat oleh keenam oksigen itu. Kanan:
[2.2.2]\omterm{def:b3:supramolecular:hosts}{kriptan}, tiga rantai dengan dua oksigen eter di antara dua kepala jembatan
nitrogen, skematis: kation dipegang di dalam sangkar tiga dimensi,
oleh enam oksigen dan kedua nitrogen.}
\end{omfigure}

\omterm{def:b3:supramolecular:hosts}{Eter mahkota} memilih kation menurut ukuran: 18-mahkota-6 mengikat $\ce{K+}$ lebih kuat daripada
$\ce{Na+}$ yang lebih kecil, yang tidak dapat mencapai keenam oksigen sekaligus, atau $\ce{Cs+}$ yang lebih besar, yang
duduk di atas cincin. \omterm{def:b3:supramolecular:hosts}{Kriptan}, yang lebih terpraorganisasi dan membungkus ion
sepenuhnya, mengikat lebih kuat lagi dan memilih lebih tajam. Tersembunyi di dalam
\omterm{def:b3:supramolecular:supramolecular}{inang}, kation itu menyeret anionnya ke dalam pelarut nonpolar, tempat anion itu,
yang tidak tersolvasi, menjadi sangat reaktif: inilah dasar katalisis transfer fase.

\section{Mengukur pengikatan}

\begin{proposition}[Pertukaran cepat]\label{prop:b3:supramolecular:fast-exchange}
Jika sebuah \omterm{def:b3:supramolecular:supramolecular}{inang} bertukar antara keadaan bebas dan keadaan terikatnya jauh lebih cepat daripada
selisih frekuensi resonansi keduanya, sinyal NMR-nya muncul pada rata-rata
$\delta_{\mathrm{obs}} = x_{\mathrm{free}}\delta_{\mathrm{free}} + x_{\mathrm{bound}}\delta_{\mathrm{bound}}$, dengan $x$ fraksi \omterm{def:b3:supramolecular:supramolecular}{inang} dalam setiap keadaan.
\end{proposition}

\begin{proof}
Setiap molekul \omterm{def:b3:supramolecular:supramolecular}{inang} mengunjungi kedua keadaan itu berkali-kali selama pengukuran;
frekuensi presesinya, yang dirata-ratakan terhadap waktu, adalah rerata berbobot keduanya, dengan
bobot berupa fraksi waktu yang dihabiskan di masing-masing keadaan, yang pada kesetimbangan sama dengan
fraksi molekul dalam setiap keadaan (\cref{ch:b3:advanced-nmr}).
\end{proof}

\begin{theorem}[Isoterm 1:1 eksak]\label{thm:b3:supramolecular:isotherm}
Untuk $\mathrm{H + G \rightleftharpoons HG}$ dengan konsentrasi total $H_0$ dan $G_0$ dan tetapan asosiasi $K$,
\[
  [\mathrm{HG}] = \frac{1}{2}\Bigl[\Bigl(H_0 + G_0 + \frac1K\Bigr) - \sqrt{\Bigl(H_0 + G_0 + \frac1K\Bigr)^2 - 4H_0G_0}\,\Bigr].
\]
\end{theorem}

\begin{proof}
Dengan $c = [\mathrm{HG}]$, $K = c/((H_0 - c)(G_0 - c))$, yaitu $c^2 - (H_0 + G_0 + 1/K)c + H_0G_0 = 0$. Dari kedua akarnya,
hanya akar yang lebih kecil yang berada di bawah $H_0$ maupun $G_0$ (hasil kali keduanya $H_0G_0$ dan jumlahnya
melebihi $H_0 + G_0$): itulah akar dengan tanda minus.
\end{proof}

\begin{omfigure}
\begin{tikzpicture}
\begin{axis}[name=a, width=0.48\linewidth, height=5.5cm, xmin=0, xmax=6, ymin=0, ymax=1.05,
  axis lines=left, xlabel={ekuivalen tamu, $G_0/H_0$}, ylabel={fraksi inang yang terikat},
  tick label style={font=\scriptsize}, label style={font=\small},
  ]
\addplot[omMeth, thick] table[x=equiv, y=K4] {figdata/out/bachelor-3-binding-titration.dat};
\addplot[omDef, thick] table[x=equiv, y=K3] {figdata/out/bachelor-3-binding-titration.dat};
\addplot[omThm, thick] table[x=equiv, y=K2] {figdata/out/bachelor-3-binding-titration.dat};
\node[font=\scriptsize, omMeth, anchor=north] at (axis cs:4,0.94) {$K = 10^4$};
\node[font=\scriptsize, omDef, anchor=north] at (axis cs:4,0.68) {$K = 10^3$};
\node[font=\scriptsize, omThm, anchor=north] at (axis cs:4,0.27) {$K = 10^2$};
\end{axis}
\begin{axis}[at={(a.east)}, anchor=west, xshift=1.4cm, width=0.48\linewidth, height=5.5cm,
  xmin=0, xmax=1, ymin=0, ymax=1.05, axis lines=left,
  xlabel={fraksi mol inang}, ylabel={$[\text{kompleks}]/c_{\text{total}}$},
  tick label style={font=\scriptsize}, label style={font=\small}]
\addplot[omDef, thick] table[x=x, y=HG] {figdata/out/bachelor-3-binding-titration-b.dat};
\addplot[omThm, thick] table[x=x, y=HG2] {figdata/out/bachelor-3-binding-titration-b.dat};
\draw[black!35, dashed] (axis cs:0.5,0) -- (axis cs:0.5,0.55) (axis cs:0.3333,0) -- (axis cs:0.3333,0.38);
\node[font=\scriptsize, omDef, anchor=south] at (axis cs:0.5,0.52) {HG};
\node[font=\scriptsize, omThm, anchor=south] at (axis cs:0.25,0.36) {$\mathrm{HG_2}$};
\end{axis}
\end{tikzpicture}

{\small Kiri: fraksi \omterm{def:b3:supramolecular:supramolecular}{inang} yang terikat selama titrasi pada $H_0 = \qty{1.0}{mmol/L}$ untuk tiga
tetapan asosiasi (dalam $\unit{L.mol^{-1}}$): makin besar $K$, makin tajam patahan pada satu ekuivalen;
jika $KH_0 \gg 1$, kurva itu hanya mengatakan bahwa $K$ besar. Kanan: \omterm{def:b3:supramolecular:job}{plot Job}, yaitu
konsentrasi kompleks pada konsentrasi total \omterm{def:b3:supramolecular:supramolecular}{inang} ditambah \omterm{def:b3:supramolecular:supramolecular}{tamu} yang tetap,
yang memuncak pada fraksi \omterm{def:b3:supramolecular:supramolecular}{inang} 1/2 untuk kompleks 1:1 dan 1/3 untuk kompleks 1:2 (tetapan
model).}
\end{omfigure}

\begin{method}[Tetapan pengikatan dari titrasi NMR]\label{met:b3:supramolecular:nmr-titration}
\begin{enumerate}
  \item Pertahankan konsentrasi \omterm{def:b3:supramolecular:supramolecular}{inang} $H_0$ tetap dan tambahkan \omterm{def:b3:supramolecular:supramolecular}{tamu}, dari 0 sampai beberapa
        ekuivalen; rekamlah sebuah sinyal \omterm{def:b3:supramolecular:supramolecular}{inang} di setiap titik (pertukaran cepat).
  \item Model: $\delta = \delta_{\mathrm{free}} + (\delta_{\mathrm{bound}} - \delta_{\mathrm{free}})[\mathrm{HG}]/H_0$, dengan $[\mathrm{HG}]$ dari isoterm eksak.
  \item Cocokkan $K$ dan $\delta_{\mathrm{bound}}$ dengan kuadrat terkecil nonlinear, dengan meminimalkan jumlah kuadrat
        sisaan; ambillah galat baku dari kelengkungan jumlah itu
        (matriks kovarians).
  \item Periksalah apakah sisaan memiliki kecenderungan (stoikiometri yang salah) dan bahwa
        fraksi terikat mencakup kira-kira 20 sampai 80\,\% (data lalu membatasi $K$).
\end{enumerate}
\end{method}

\begin{definition}[Plot Job]\label{def:b3:supramolecular:job}
Yang disebut \emph{plot Job}\index{plot Job} (metode variasi kontinu) adalah alur
konsentrasi sebuah kompleks, atau sinyal yang sebanding dengannya, terhadap
fraksi mol salah satu komponen, pada konsentrasi total keduanya yang
tetap.
\end{definition}

\begin{proposition}[Maksimum plot Job]\label{prop:b3:supramolecular:job}
Untuk satu kompleks tunggal $\mathrm{H_nG_m}$, \omterm{def:b3:supramolecular:job}{plot Job} mencapai maksimum pada fraksi mol \omterm{def:b3:supramolecular:supramolecular}{inang}
$x = n/(n + m)$, berapa pun tetapan pengikatannya.
\end{proposition}

\begin{proof}
Dengan konsentrasi total $T$, $H_t = xT$, $G_t = (1 - x)T$, kompleks $c$, bebas $[\mathrm H] = H_t - nc$,
$[\mathrm G] = G_t - mc$, dan $c = \beta[\mathrm H]^n[\mathrm G]^m$. Dengan menurunkan terhadap $x$ pada maksimum, tempat
$\dd c/\dd x = 0$: $0 = \beta\bigl(n[\mathrm H]^{n-1}[\mathrm G]^m T - m[\mathrm H]^n[\mathrm G]^{m-1}T\bigr)$, sehingga $n[\mathrm G] = m[\mathrm H]$. Dengan menyubstitusikannya,
$n(G_t - mc) = m(H_t - nc)$, sehingga $nG_t = mH_t$, yaitu $n(1 - x) = mx$: $x = n/(n + m)$.
\end{proof}

\begin{method}[Menjalankan plot Job]\label{met:b3:supramolecular:job}
\begin{enumerate}
  \item Siapkan larutan stok \omterm{def:b3:supramolecular:supramolecular}{inang} dan \omterm{def:b3:supramolecular:supramolecular}{tamu} yang ekuimolar.
  \item Campurkan keduanya dengan proporsi $x$ dari 0 sampai 1 pada volume total yang tetap.
  \item Ukurlah sebuah sifat yang sebanding dengan kompleks (absorbansi yang dikoreksi
        untuk spesies bebas, atau $x\,\Delta\delta$ dalam NMR).
  \item Bacalah stoikiometrinya dari posisi maksimum.
\end{enumerate}
\end{method}

\begin{definition}[Kalorimetri titrasi isotermal]\label{def:b3:supramolecular:itc}
Yang disebut \emph{kalorimetri titrasi isotermal}\index{kalorimetri titrasi isotermal}
(ITC) adalah teknik yang mengukur kalor yang dilepaskan atau diserap pada setiap injeksi larutan \omterm{def:b3:supramolecular:supramolecular}{tamu}
ke dalam larutan \omterm{def:b3:supramolecular:supramolecular}{inang} yang dijaga pada suhu tetap.
\end{definition}

\begin{proposition}[Kalor per injeksi]\label{prop:b3:supramolecular:itc}
Dengan mengabaikan pengenceran, kalor injeksi ke-$i$ ke dalam sel bervolume $V_0$ adalah
$q_i = \Delta H^\circ V_0\bigl([\mathrm{HG}]_i - [\mathrm{HG}]_{i-1}\bigr)$.
\end{proposition}

\begin{proof}
Kalor itu adalah entalpi reaksi kali jumlah kompleks yang terbentuk selama
injeksi, dan jumlah itu adalah perubahan $[\mathrm{HG}]$ kali volume sel.
\end{proof}

\begin{method}[Membaca isoterm ITC]\label{met:b3:supramolecular:itc}
\begin{enumerate}
  \item Alurkan kalor per mol \omterm{def:b3:supramolecular:supramolecular}{tamu} yang diinjeksikan terhadap rasio molar.
  \item Tinggi dataran sebelum kejenuhan memberikan $\Delta H^\circ$; rasio molar pada
        titik belok memberikan stoikiometri; kecuraman memberikan $K$.
  \item Cocokkan ketiganya dengan isoterm eksak; lalu $\Delta G^\circ = -RT\ln K$ dan
        $T\Delta S^\circ = \Delta H^\circ - \Delta G^\circ$.
\end{enumerate}
\end{method}

Satu percobaan ITC dengan demikian memberikan seluruh tanda termodinamika
pengikatan, entalpi dan entropi, yang hanya dapat diberikan oleh titrasi NMR atau UV--tampak
melalui serangkaian suhu.

\begin{inthelab}[Sebuah titrasi NMR]
Larutan \omterm{def:b3:supramolecular:supramolecular}{inang} dalam pelarut terdeuterasi, sekitar satu milimol per liter,
dimasukkan ke dalam tabung NMR dan spektrumnya direkam. Larutan pekat
\omterm{def:b3:supramolecular:supramolecular}{tamu} yang dibuat di dalam larutan \omterm{def:b3:supramolecular:supramolecular}{inang}, agar konsentrasi \omterm{def:b3:supramolecular:supramolecular}{inang} tetap
konstan, ditambahkan sedikit demi sedikit dengan mikrosuntik; setelah setiap penambahan,
tabung dikocok, disetimbangkan pada suhu probe, dan spektrumnya
direkam. Geseran sebuah sinyal \omterm{def:b3:supramolecular:supramolecular}{inang} yang bergerak dibaca di setiap titik lalu
dicocokkan.
\end{inthelab}

\section{Inang}

\omterm{def:b3:supramolecular:hosts}{Siklodekstrin}, yang dibuat oleh \omterm{def:b3:complex-kinetics:enzyme}{enzim} dari pati, memiliki enam, tujuh, atau delapan satuan
glukosa ($\alpha$, $\beta$, $\gamma$). Rongga hidrofobiknya menampung bagian-bagian nonpolar \omterm{def:b3:supramolecular:supramolecular}{tamu} di dalam
air: \omterm{def:b3:supramolecular:hosts}{siklodekstrin} melarutkan obat, membawa wewangian, dan menyingkirkan kolesterol dari
membran sel. \omterm{def:b3:supramolecular:hosts}{Kaliksarena} dan kukurbituril yang berbentuk labu, yaitu
makrosiklus kaku dari satuan-satuan glikoluril dengan gerbang yang berjajar karbonil, mengikat kation dan
kation organik dengan sangat kuat di dalam air.

\begin{omfigure}
\begin{tikzpicture}[font=\small]
  \draw[black!70, fill=omDef!8] (-1.9,1.2) -- (-1.2,-1.2) arc[start angle=180, end angle=360, x radius=1.2, y radius=0.3]
    -- (1.9,1.2);
  \draw[black!70, fill=omDef!18] (0,1.2) ellipse[x radius=1.9, y radius=0.45];
  \draw[black!40, dashed] (-1.2,-1.2) arc[start angle=180, end angle=0, x radius=1.2, y radius=0.3];
  \node[draw=omThm, thick, regular polygon, regular polygon sides=6, minimum size=8mm, inner sep=0pt] at (0,0.2) {};
  \draw[omThm, thick] (0,0.62) -- (0,2.0);
  \node[font=\scriptsize, omThm, anchor=west] at (0.1,2.0) {tamu};
  \node[font=\scriptsize, anchor=west] at (2.0,1.2) {bibir lebar: OH sekunder};
  \node[font=\scriptsize, anchor=west] at (1.3,-1.25) {bibir sempit: OH primer};
  \node[font=\scriptsize, anchor=east] at (-1.75,0) {rongga hidrofobik};
\end{tikzpicture}

{\small Sebuah \omterm{def:b3:supramolecular:hosts}{siklodekstrin}, yang digambar sebagai kerucut terpancung yang dibentuknya, dengan \omterm{def:b3:supramolecular:supramolecular}{tamu}
aromatik di dalam rongganya. Gugus-gugus hidroksi pada kedua bibirnya membuat bagian luar
larut dalam air; bagian dalam, yang berlapis gugus C--H dan oksigen glikosida,
bersifat nonpolar.}
\end{omfigure}

\section{Perakitan mandiri dan mesin molekul}

\begin{definition}[Perakitan mandiri]\label{def:b3:supramolecular:self-assembly}
Yang disebut \emph{perakitan mandiri}\index{perakitan mandiri} adalah penataan spontan dan reversibel
komponen-komponen menjadi struktur tertentu, di bawah kendali
informasi yang terkandung dalam komponen-komponen itu sendiri: bentuk dan situs-situs
pengikatannya.
\end{definition}

Ion logam dengan sudut koordinasi yang tetap dan ligan penjembatan yang kaku merakit diri
menjadi persegi, sangkar, dan kapsul dalam satu tahap, karena setiap ikatan dapat membuka
dan menutup sampai struktur yang paling stabil, tanpa tegangan atau situs yang menggantung,
tercapai. Pemasangan basa dalam DNA, dan pelipatan protein, adalah gagasan yang sama
dalam skala yang lebih besar.

\begin{definition}[Molekul yang saling mengunci secara mekanis]\label{def:b3:supramolecular:interlocked}
Yang disebut \emph{molekul yang saling mengunci secara mekanis}\index{molekul yang saling mengunci secara mekanis}
adalah molekul yang bagian-bagiannya tidak saling terikat tetapi tidak dapat dipisahkan
tanpa memutus sebuah ikatan kovalen. Yang disebut \emph{rotaksana}\index{rotaksana} adalah cincin
yang terulir pada sebuah poros dengan penyumbat meruah di kedua ujungnya; sedangkan
\emph{katenana}\index{katenana} terdiri atas cincin-cincin yang saling mengunci.
\end{definition}

\begin{definition}[Sintesis templat]\label{def:b3:supramolecular:template}
Yang disebut \emph{sintesis templat}\index{sintesis templat} adalah sintesis yang memakai interaksi
sementara, sering dengan ion logam, untuk mempertahankan komponen-komponen reaktif dalam
geometri yang diperlukan bagi reaksi yang jika tidak demikian kecil kemungkinannya, seperti
penutupan sebuah cincin di sekeliling cincin lain.
\end{definition}

\begin{omfigure}
\begin{tikzpicture}[font=\small, >=Stealth]
  % tahap 1: dua bulan sabit di sekeliling Cu+
  \draw[omDef, line width=2pt] (0.3,0) arc[start angle=0, end angle=300, radius=0.75];
  \draw[omThm, line width=2pt] (-0.3,0) arc[start angle=180, end angle=-120, radius=0.75];
  \fill[omProp] (0,0) circle (0.13);
  \node[font=\scriptsize] at (0,-1.35) {$\ce{Cu+}$ memegang dua};
  \node[font=\scriptsize] at (0,-1.7) {ligan terbuka bersilangan};
  \draw[->, thick] (1.4,0) -- (2.3,0) node[midway, above, font=\scriptsize] {tutup};
  % tahap 2: cincin-cincin tertutup yang saling mengunci dengan Cu
  \begin{scope}[xshift=3.6cm]
    \draw[omDef, line width=2pt] (-0.35,0) circle (0.75);
    \draw[omThm, line width=2pt] (0.35,0) circle (0.75);
    \draw[white, line width=5pt] ([shift={(52:0.75)}]-0.35,0) arc[start angle=52, end angle=72, radius=0.75];
    \draw[omDef, line width=2pt] ([shift={(48:0.75)}]-0.35,0) arc[start angle=48, end angle=76, radius=0.75];
    \fill[omProp] (0,0) circle (0.13);
    \node[font=\scriptsize] at (0,-1.35) {katenat};
  \end{scope}
  \draw[->, thick] (5.0,0) -- (5.9,0) node[midway, above, font=\scriptsize] {$-\ce{Cu+}$};
  \begin{scope}[xshift=7.2cm]
    \draw[omDef, line width=2pt] (-0.35,0) circle (0.75);
    \draw[omThm, line width=2pt] (0.35,0) circle (0.75);
    \draw[white, line width=5pt] ([shift={(52:0.75)}]-0.35,0) arc[start angle=52, end angle=72, radius=0.75];
    \draw[omDef, line width=2pt] ([shift={(48:0.75)}]-0.35,0) arc[start angle=48, end angle=76, radius=0.75];
    \node[font=\scriptsize] at (0,-1.35) {[2]katenana};
  \end{scope}
  % rotaksana
  \begin{scope}[xshift=10.2cm]
    \draw[black!70, line width=1.6pt] (-1.1,0) -- (1.1,0);
    \fill[black!60] (-1.25,0) circle (0.2) (1.25,0) circle (0.2);
    \draw[omThm, line width=2pt] (0,0) ellipse[x radius=0.18, y radius=0.55];
    \node[font=\scriptsize] at (0,-1.35) {[2]rotaksana};
  \end{scope}
\end{tikzpicture}

{\small \omterm{def:b3:supramolecular:template}{Sintesis templat} sebuah \omterm{def:b3:supramolecular:interlocked}{katenana} (skematis): ion tembaga(I) memegang dua
ligan berbasis fenantrolina yang bersilangan tegak lurus dalam koordinasi
tetrahedralnya; menutup masing-masing menjadi cincin memberikan dua cincin yang saling mengunci di sekeliling
logam, dan penyingkiran tembaga menyisakan \omterm{def:b3:supramolecular:interlocked}{katenana} bebas. Kanan: sebuah \omterm{def:b3:supramolecular:interlocked}{rotaksana}, yaitu
cincin yang terperangkap pada poros oleh dua penyumbat.}
\end{omfigure}

\begin{definition}[Mesin molekul]\label{def:b3:supramolecular:machine}
Yang disebut \emph{mesin molekul}\index{mesin molekul} adalah rakitan yang
komponen-komponennya bergerak relatif satu sama lain secara terkendali sebagai tanggapan atas
sebuah rangsangan (cahaya, perubahan redoks, perubahan pH), dan dapat dibuat untuk melakukan kerja.
\end{definition}

Dalam ulang-alik molekul, sebuah cincin pada poros \omterm{def:b3:supramolecular:interlocked}{rotaksana} menempati salah satu dari dua stasiun;
rangsangan yang melemahkan ikatannya dengan stasiun pertama mengirimnya ke stasiun kedua, dan
rangsangan kebalikannya mengirimnya kembali. Motor putar yang digerakkan cahaya berbasis
alkena yang sesak berputar hanya ke satu arah, dengan berselang-selingnya isomerisasi
fotokimia dan tahap-tahap termal yang tidak dapat berjalan mundur. Laju
ulang-alik di antara stasiun-stasiun diukur dengan NMR suhu variabel, dari
koalesensi sinyal kedua bentuk (\cref{ch:b3:advanced-nmr}).

\begin{safety}{\ghs[0.62cm]{GHS07}}
% ledger: ghs4:18C6
18-Mahkota-6: berbahaya jika tertelan, mengiritasi kulit, mata, dan saluran pernapasan.
\omterm{def:b3:supramolecular:hosts}{Eter mahkota} membawa garam logam ke dalam pelarut organik dan menembus kulit; zat itu
ditimbang dan ditangani dengan sarung tangan.
\end{safety}

\begin{history}[Kimia supramolekul dan mesin-mesinnya]
\omterm{def:b3:supramolecular:hosts}{Eter mahkota} Charles Pedersen (1967), \omterm{def:b3:supramolecular:hosts}{kriptan} Jean-Marie Lehn, dan \omterm{def:b3:supramolecular:supramolecular}{inang} terpraorganisasi Donald
Cram mendirikan bidang ini; ketiganya berbagi Hadiah Nobel Kimia
1987. Jean-Pierre Sauvage membuat \omterm{def:b3:supramolecular:interlocked}{katenana} dengan rendemen tinggi melalui
templat tembaga pada 1983, Fraser Stoddart membangun ulang-alik \omterm{def:b3:supramolecular:interlocked}{rotaksana}, dan Ben
Feringa membuat motor putar yang digerakkan cahaya pada 1999; mereka berbagi Hadiah Nobel Kimia 2016
untuk perancangan dan sintesis \omterm{def:b3:supramolecular:machine}{mesin molekul}.
\end{history}

\section{Latihan}

\begin{exercise}[$\star$]\label{exo:b3:supramolecular:1}
Sebutkan interaksi utama dalam: $\ce{K+}$ dalam 18-mahkota-6; dua cincin benzena yang saling berhadapan muka;
$\ce{K+}$ di atas cincin benzena; iodopentafluorobenzena dengan piridina; pasangan
guanina--sitosina.
\end{exercise}

\begin{exercise}[$\star$]\label{exo:b3:supramolecular:2}
Hitunglah $\Delta G^\circ$ pada \qty{298}{K} untuk \omterm{def:b3:supramolecular:supramolecular}{kompleks inang--tamu} dengan $K_a = \qty{1.0e4}{L.mol^{-1}}$.
\end{exercise}

\begin{exercise}[$\star$]\label{exo:b3:supramolecular:3}
Dalam metanol, $\log K = 6.1$ untuk $\ce{K+}$ dan 4.4 untuk $\ce{Na+}$ dengan 18-mahkota-6 (data latihan).
Hitunglah selektivitasnya dan selisih $\Delta G^\circ$-nya.
\end{exercise}

\begin{exercise}[$\star$]\label{exo:b3:supramolecular:4}
Sebuah \omterm{def:b3:supramolecular:job}{plot Job} memuncak pada fraksi mol \omterm{def:b3:supramolecular:supramolecular}{inang} 0.33. Berapakah stoikiometrinya?
\end{exercise}

\begin{exercise}[$\star\star$]\label{exo:b3:supramolecular:5}
Sebuah \omterm{def:b3:supramolecular:supramolecular}{inang} pada \qty{2.0}{mmol/L} dengan $K = \qty{1500}{L.mol^{-1}}$, $\delta_{\mathrm{free}} = \qty{3.560}{ppm}$, dan $\delta_{\mathrm{bound}} = \qty{3.720}{ppm}$ menerima satu
ekuivalen \omterm{def:b3:supramolecular:supramolecular}{tamu}. Hitunglah fraksi terikat dan geseran yang teramati.
\end{exercise}

\begin{exercise}[$\star\star$]\label{exo:b3:supramolecular:6}
Untuk nikel(II), $\log\beta_6 = 8.6$ dengan amonia dan $\log\beta_3 = 18.3$ dengan etana-1,2-diamina (data
latihan). Hitunglah tetapan dan $\Delta G^\circ$ pada \qty{298}{K} untuk
$\ce{[Ni(NH3)6]^2+} + 3\,\mathrm{en} \rightleftharpoons \ce{[Ni(en)3]^2+} + \ce{6NH3}$, dan sebutkan apa yang mendorongnya.
\end{exercise}

\begin{exercise}[$\star\star$]\label{exo:b3:supramolecular:7}
Sebuah percobaan ITC pada \qty{298}{K} memberikan $K = \qty{2.0e5}{L.mol^{-1}}$ dan $\Delta H^\circ = \qty{-40}{kJ/mol}$. Hitunglah $\Delta G^\circ$, $T\Delta S^\circ$, dan
$\Delta S^\circ$.
\end{exercise}

\begin{exercise}[$\star\star$]\label{exo:b3:supramolecular:8}
Jelaskan peran tembaga(I), dan geometri tetrahedralnya, dalam sintesis \omterm{def:b3:supramolecular:interlocked}{katenana}
Sauvage. Bagaimana tembaga itu disingkirkan pada akhirnya?
\end{exercise}

\begin{exercise}[$\star\star$]\label{exo:b3:supramolecular:9}
Kedua sinyal sebuah ulang-alik \omterm{def:b3:supramolecular:interlocked}{rotaksana}, yang terpisah \qty{120}{Hz}, berkoalesensi pada \qty{300}{K} (data
latihan). Hitunglah tetapan laju pertukaran pada koalesensi dan \omterm{def:b3:rate-theories:activation}{energi Gibbs
aktivasi} ulang-alik itu.
\end{exercise}

\begin{exercise}[$\star\star\star$]\label{exo:b3:supramolecular:10}
Sebuah obat memiliki kelarutan intrinsik \qty{0.10}{mmol/L} dan membentuk kompleks 1:1 dengan sebuah
\omterm{def:b3:supramolecular:hosts}{siklodekstrin}, $K = \qty{500}{L.mol^{-1}}$ (data latihan). Tunjukkan bahwa kelarutan total dalam
larutan \omterm{def:b3:supramolecular:hosts}{siklodekstrin} total $C_t$ adalah $S_0 + KS_0C_t/(1 + KS_0)$, dan hitunglah nilainya untuk $C_t = \qty{10}{mmol/L}$.
\end{exercise}

\begin{exercise}[$\star\star\star$]\label{exo:b3:supramolecular:11}
Tunjukkan dari isoterm eksak bahwa jika $KH_0 \gg 1$, fraksi terikat sama dengan
jumlah ekuivalen sampai satu, lalu tetap satu. Mengapa titrasi semacam itu
hanya memberikan batas bawah untuk $K$?
\end{exercise}

\begin{exercise}[$\star\star\star$]\label{exo:b3:supramolecular:12}
Untuk \omterm{def:b3:supramolecular:supramolecular}{inang} pada \cref{exo:b3:supramolecular:5}, anggaplah asosiasinya
\omterm{def:b3:rate-theories:diffusion}{dikendalikan difusi}, $k_{\mathrm{on}} = \qty{1e9}{L.mol^{-1}.s^{-1}}$. Hitunglah $k_{\mathrm{off}}$ dan tunjukkan bahwa pertukarannya cepat dalam
skala waktu NMR pada \qty{400}{MHz}.
\end{exercise}

\section{Soal: Benzena Ungu}

\begin{problem}[{Soal akhir pekan --- benzena ungu: kalium dalam eter mahkota,
ekstraksi permanganat ke dalam pelarut organik, tetapan pengikatan
dari titrasi NMR, dan mengapa sejumlah katalitik eter mahkota sudah cukup}]
\label{pb:b3:supramolecular:1}
% ledger: const:R, ghs:KMnO4, ghs:toluene, rion:K+VI, rion:Na+VI
Data soal. Dalam metanol pada \qty{298}{K}, $\log K = 6.1$ untuk $\ce{K+}$ dan 4.4 untuk $\ce{Na+}$ dengan
18-mahkota-6. Ekstraksi: fase air dengan \qty{0.10}{mol/L} $\ce{KCl}$ dan \qty{1.0}{mmol/L} $\ce{KMnO4}$ dikocok
dengan toluena bervolume sama yang mengandung \omterm{def:b3:supramolecular:hosts}{eter mahkota} L; satu-satunya spesies
yang terekstraksi adalah pasangan ion $\ce{[KL]+}\ce{MnO4-}$, dengan
$K_{\mathrm{ex}} = [\ce{[KL]+MnO4-}]_{\mathrm{org}}/([\ce{K+}]_{\mathrm{aq}}[\ce{MnO4-}]_{\mathrm{aq}}[\mathrm L]_{\mathrm{org}}) = \qty{1.0e5}{L^2.mol^{-2}}$, dan \omterm{def:b3:supramolecular:hosts}{eter mahkota} tetap berada di dalam toluena. Titrasi NMR sebuah
\omterm{def:b3:supramolecular:hosts}{eter mahkota} (\qty{2.0}{mmol/L}) dengan garam natrium dalam metanol terdeuterasi, satu sinyal $\ce{CH2}$
(ppm) terhadap ekuivalen $\ce{Na+}$:

\begin{center}\small
\begin{tabular}{l*{11}{c}}
\hline
ekuiv. & 0 & 0.25 & 0.5 & 0.75 & 1.0 & 1.5 & 2.0 & 3.0 & 4.0 & 6.0 & 8.0\\
$\delta$ & 3.560 & 3.590 & 3.612 & 3.635 & 3.649 & 3.674 & 3.688 & 3.697 & 3.704 & 3.708 & 3.714\\
\hline
\end{tabular}
\end{center}
% tabel: binding-titration-c

\medskip\noindent\textbf{Bagian I --- Kompleks.}
\begin{enumerate}[beginpenalty=10000]
  \item Gambarlah 18-mahkota-6 dan kompleks kaliumnya.
  \item Hitunglah $\Delta G^\circ$ pengikatan $\ce{K+}$ dan $\ce{Na+}$.
  \item Hitunglah selektivitas $\ce{K+}/\ce{Na+}$.
  \item Dengan jari-jari berkoordinasi enam $r(\ce{K+}) = \qty{138}{pm}$ dan $r(\ce{Na+}) = \qty{102}{pm}$, jelaskan selektivitas itu.
  \item Mengapa kompleks itu larut dalam toluena sedangkan $\ce{KCl}$ tidak?
  \item Mengapa \omterm{def:b3:supramolecular:hosts}{eter mahkota} mengikat lebih lemah di dalam air daripada di dalam metanol?
\end{enumerate}

\medskip\noindent\textbf{Bagian II --- Ekstraksi.}
\begin{enumerate}[resume, beginpenalty=10000]
  \item Tuliskan kesetimbangan ekstraksinya.
  \item Dengan $y$ konsentrasi yang terekstraksi, tuliskan persamaan untuk $y$ jika
        $[\mathrm L]_0 = \qty{1.0}{mmol/L}$.
  \item Selesaikan persamaan itu, dan berikan fraksi permanganat yang terekstraksi.
  \item Hitunglah fraksi itu dengan $[\mathrm L]_0 = \qty{10}{mmol/L}$.
  \item Hitunglah fraksi itu dengan $[\mathrm L]_0 = \qty{0.10}{mmol/L}$.
  \item Mengapa toluena itu berwarna ungu, dan mengapa permanganat yang terekstraksi merupakan oksidator
        yang lebih kuat daripada di dalam air?
\end{enumerate}

\medskip\noindent\textbf{Bagian III --- Titrasi NMR.}
\begin{enumerate}[resume, beginpenalty=10000]
  \item Mengapa satu sinyal rata-rata yang bergeser, dan bukan dua sinyal?
  \item Tuliskan model untuk $\delta$ terhadap konsentrasi \omterm{def:b3:supramolecular:supramolecular}{tamu}.
  \item Cocokkan $K$ dan $\delta_{\mathrm{bound}}$ dengan kuadrat terkecil; berikan $K$ beserta galat bakunya.
  \item Berikan sebaran sisaannya dan berilah komentar.
  \item Dalam rentang ekuivalen manakah fraksi terikat berada di antara 20 dan 80\,\%?
  \item Mengapa titrasi yang sama akan gagal mengukur $K$ untuk kalium, $\log K = 6.1$?
\end{enumerate}

\medskip\noindent\textbf{Bagian IV --- Katalisis transfer fase.}
\begin{enumerate}[resume, beginpenalty=10000]
  \item Sebuah alkena dalam toluena dioksidasi oleh permanganat yang terekstraksi. Di manakah
        reaksi itu berlangsung?
  \item Apakah yang terjadi pada \omterm{def:b3:supramolecular:hosts}{eter mahkota} setelah reaksi, dan mengapa sejumlah katalitik
        sudah cukup?
  \item Mengapa $\ce{KCl}$ berlebih di dalam air berguna?
  \item Katalis-katalis lebih murah apakah yang melakukan tugas yang sama dalam industri?
  \item Mengapa sekarang dipakai toluena, sedangkan Pedersen memakai benzena?
  \item Nyatakan hasilnya: fraksi permanganat yang terekstraksi pada
        $[\mathrm L]_0 = \qty{1.0}{mmol/L}$.
\end{enumerate}
\end{problem}
